%% sections/04-results.tex

\section{Results}
\label{sec:results}

\subsection{Innovation Rate by Session and Campaign}

Table~\ref{tab:innovations} reports innovations logged per session across the 7 campaigns.

\begin{table}[t]
\caption{Innovations Logged Per Campaign: Session Range, Total, Mean per Session}
\label{tab:innovations}
\begin{tabular}{@{}lllrrl@{}}
\toprule
Campaign & Rubric at open & Rubric at close & Total innov. & Mean/session & Campaign archetype \\
\midrule
C1 & v1.0 & v1.4 & 40+ & 5.7 & grief-themed \\
C2 & v1.4 & v1.6 & 12 & 1.7 & covenant-restoration \\
C3 & v1.6 & v1.8 & 16 & 2.3 & faith/authority \\
C4 & v1.8 & v1.12 & 16 & 2.3 & siege-resource \\
C5 & v1.12 & v1.12 & 8 & 1.1 & professional-code \\
C6 & v1.12 & v1.12 & 15 & 2.1 & duty-spine \\
C7 & v1.12 & v1.12 & 7 & 1.0 & no-designed-stakes \\
\bottomrule
\end{tabular}
\end{table}

The self-calibration pattern is visible: C1 has the highest mean innovation rate (5.7
per session) because the rubric begins at v1.0 with no accumulated anchors. C2 drops
sharply (1.7 per session) because the C1 amendments absorbed the most common quality
patterns. C3 and C4 produce moderate bursts (2.3 per session each) as new campaign
archetypes introduce quality phenomena the v1.4 and v1.6 anchors had not captured.
C5--C7, all running under the mature v1.12 rubric, produce the lowest innovation rates
(1.0--2.1 per session).

It is notable that C6 (Military Unit, duty-spine) produces 15 innovations despite
running under the mature v1.12 rubric. Inspection of the C6 innovation log reveals
that the Unit Cohesion Die mechanic produced several quality phenomena not seen in
any prior campaign: ``grief-paragraph resource expression'' (a PC's resource decisions
shaped by their grief-paragraph commitment rather than tactical calculation),
``unit-member-as-named-loss'' atmospheric landing, and ``order-defiance as character
commitment'' agency pattern. These were all universal-scope innovations (likely to
recur in any duty-spine campaign with a Cohesion Die), but they required the
combination of duty-spine archetype and Cohesion Die mechanic --- neither had appeared
alone in the corpus.

\subsection{Cluster Coherence}

Table~\ref{tab:cluster-coherence} characterizes the thematic coherence of amendment-triggering
clusters by reporting the number of distinct quality phenomena addressed and the degree
of cross-session repetition.

\begin{table}[t]
\caption{Cluster Coherence: All 15 Amendment-Triggering Dimension-Clusters}
\label{tab:cluster-coherence}
\begin{tabular}{@{}llrcl@{}}
\toprule
Version & Dimension & Cluster size & Sessions spanned & Thematic label \\
\midrule
v1.1 & MF & 2 & S01 & Cross-dimensional damage \\
v1.1 & AL & 3 & S01 & Per-PC reception + chain-reaction \\
v1.2 & MFid & 4 & S01--S02 & Anchor-level single-gate failure \\
v1.2 & AL & 3 & S02 & Finder/receiver divergence + silent reception \\
v1.3 & MFid & 3 & S03 & Positional-geometry failure \\
v1.4 & MF & 2 & S03--S04 & Cross-character compound \\
v1.4 & AL & 2 & S04--S05 & Multi-scene silent-reception chain \\
v1.5 & AL & 3 & S07, C2-S01--S02 & Cross-session act-without-announcement \\
v1.6 & AL & 2 & S07, C2-S04 & Simultaneous arc-convergence \\
v1.7 & AL & 2 & C3-S01 & Framework-failure arc-activation \\
v1.8 & AL & 2 & C3-S03--S04 & Arc-activation through behavioral stance \\
v1.9 & AL, CA & 4 & C4-S01 & Grief-paragraph resonance + sustained pattern \\
v1.10 & CA & 2 & C4-S02--S03 & Character-register conversion \\
v1.11 & MF & 2 & C4-S03--S04 & Portent-as-directed-foreknowledge \\
v1.12 & AL & 2 & C4-S02--S03 & Intra-session restraint \\
\bottomrule
\end{tabular}
\end{table}

All 15 clusters showed thematic unity: innovations within a cluster addressed the same
underlying quality phenomenon from different session contexts. No cluster contained
innovations addressing genuinely distinct phenomena that happened to share a dimension
tag. In particular, the AL clusters were consistently distinguishable from each other:
v1.1 (who receives which beat), v1.2 (finder and receiver separated), v1.4 (reception
persisting across scenes), v1.5 (reception expressing across sessions), v1.6 (reception
releasing into simultaneous convergence), v1.7 (reception from framework-failure rather
than discovery), v1.8 (reception from behavioral stance rather than speech), v1.9
(reception via grief-pattern recognition), v1.12 (reception expressing intra-session)
are nine distinct quality phenomena, all in the Atmospheric Landing dimension, each
clearly distinguishable from the others.

\subsection{Innovation Taxonomy}

Figure~\ref{fig:taxonomy} shows the distribution of innovations across scope
classifications.

\begin{figure}[t]
  \centering
  \fbox{\parbox{0.85\columnwidth}{\centering
    [Bar chart: innovations by scope classification. \\
    Universal: 78 (65\%). \\
    Adventure-specific: 18 (15\%). \\
    Party-specific: 24 (20\%). \\
    Total: 120 entries analyzed. \\
    All 15 amendment clusters contained at least 1 Universal-scope innovation.]}}
  \caption{Innovation scope distribution. Universal-scope innovations are the
    primary raw material for amendments. Party-specific and adventure-specific
    innovations serve as corroborating evidence only.}
  \label{fig:taxonomy}
\end{figure}

Of 120+ logged innovations, approximately 78 (65\%) were classified as Universal.
The remaining 35\% were adventure-specific (15\%) or party-specific (20\%). All 15
amendment-triggering clusters contained at least one Universal-scope innovation;
no amendment was adopted on party-specific-only evidence.

\subsection{Amendment Rate Decay and Archetype Reactivation}

Figure~\ref{fig:amendment-rate} shows the amendment adoption rate (amendments per session)
across the 49-session corpus.

\begin{figure}[t]
  \centering
  \fbox{\parbox{0.85\columnwidth}{\centering
    [Line chart: amendments adopted per session window. \\
    C1 (S01-S07): 0.57 per session (4 adopted). \\
    C2 (C2-S01 to C2-S07): 0.29 per session (2 adopted). \\
    C3 (C3-S1 to C3-S7): 0.29 per session (2 adopted). \\
    C4 (C4-S1 to C4-S7): 0.57 per session (4 adopted). \\
    C5--C7 (S1-S7 each): 0.05 per session total (0 adopted after v1.12). \\
    Dashed line: 0.25 per session (approximate steady-state prediction).]}}
  \caption{Amendment adoption rate by session position. C1 and C4 show high rates
    (novel rubric and novel archetype, respectively). C5--C7 approach zero
    as the v1.12 rubric stabilizes.}
  \label{fig:amendment-rate}
\end{figure}

The decay-and-reactivation pattern is clear. C1's rate (0.57 per session) reflects a
fresh rubric encountering 7 sessions of quality outcomes for the first time. C2's rate
drops to 0.29 as the rubric absorbs the most common patterns --- per-PC reception
tracking, cross-character compound mechanics, and anchor-level fallback requirements
were the highest-frequency gaps in the v1.0 rubric. C3 holds at 0.29 as the
faith/authority archetype introduces 2 new anchor clusters (framework-failure
arc-activation from a fumbled Insight roll; behavioral-stance arc-activation from
combat-as-precondition). C4 reactivates at 0.57 as the siege-resource mechanic produces
4 new clusters the grief-adjacent and covenant anchors had not addressed.

C5 through C7, all under v1.12, produce no adoptions. This does not mean the rubric
is complete for all future campaign archetypes: it means the three archetypes run
under v1.12 (professional-code, duty-spine, no-stakes) produced innovations that were
either already captured by v1.12 anchors or were adventure-specific and party-specific
--- not universal enough to trigger amendments.

\subsection{The Atmospheric Landing Asymmetry in Detail}

Eight of the 13 amendments (and 9 of the 15 amendment dimension-clusters) affected
Atmospheric Landing. This asymmetry deserves analysis beyond noting that AL is complex.

A dimension-by-dimension breakdown of why AL produced more clusters than others:

\begin{description}
  \item[Engagement (0 amendments):] Engagement measures PC participation and collective
    unscripted outcomes. These properties have a bounded quality space: either PCs drove
    scenes or they did not; either the outcome was unscripted or it was not. The initial
    anchors were adequate.

  \item[Mechanical Fairness (3 amendments):] MF produced 3 amendments (v1.1, v1.4, v1.11),
    all focused on the 10-band's definition of what constitutes the most sophisticated
    mechanical play. The cross-dimensional-harmlessness anchor (v1.1) was needed early;
    the cross-character compound anchors (v1.4, v1.11) required 2 sessions of Portent
    evidence to identify.

  \item[Pacing (0 amendments):] Pacing is architectural: did scenes open and close cleanly?
    This is assessable from the session log without fine-grained anchor specification.

  \item[Character Agency (2 amendments):] CA produced 2 amendments (v1.9, v1.10), both
    expanding the 9--10 band to include behavioral-pattern forms of agency that the
    initial anchors did not describe.

  \item[Module Fidelity (2 amendments):] MFid produced 2 amendments (v1.2, v1.3),
    both addressing anchor-level symptom reliability --- the most practically important
    MFid gaps revealed by early sessions.

  \item[Atmospheric Landing (8 amendments):] AL spans the entire space of how emotional
    and atmospheric intent can be realized in session play. This space is effectively
    unbounded, as evidenced by the distinct phenomena addressed in each AL amendment.
    The initial anchors (reception tracking, inter-PC chains) were necessary foundations;
    each amendment added a distinct realization mode.

  \item[Surprise (0 amendments):] Surprise is scored by tag-counting (7+ tags at 10-band)
    and does not require fine-grained anchor specification.

  \item[Table Readiness (0 amendments):] Table Readiness is assessable from a binary
    check (did the DM need outside reference?) without fine-grained anchor specification.
\end{description}

The Atmospheric Landing asymmetry is structurally expected: the more a dimension
describes complex, emergent, inter-PC phenomena, the more likely it is to produce
innovation clusters as new phenomena are observed.
